\section{Methods}
\label{sec:methods}

\subsection{Data and effect vectors}
\label{sec:methods-data}
The primary dictionary is the genome-scale CRISPRi Perturb-seq screen of
\citet{Zhu2025} in primary human CD4\textsuperscript{+} T cells (Marson and Pritchard
laboratories, CZI Virtual Cells Platform), comprising differential-expression statistics for
$33{,}983$ single-gene knockdowns across $10{,}282$ measured genes from four donors. For each
knockdown $j$ we take the published per-gene effect vector $\Emat_{:,j}\in\mathbb{R}^{G}$,
the signed transcriptional shift the knockdown induces relative to non-targeting controls.
Stacking these columns gives the effect matrix $\Emat\in\mathbb{R}^{G\times P}$. A
\emph{target} is a desired cell-state shift $\dvec\in\mathbb{R}^{G}$, computed as the
difference between the mean expression profile of the source state and that of the
destination state (e.g. T\textsubscript{H}2 minus T\textsubscript{H}1). All effect and target
vectors are restricted to the genes measured in both the screen and the target contrast, and
comparisons are reported as cosine similarity so that magnitude scaling does not affect the
verdict. The identical pipeline is applied without modification to the K562 combinatorial
CRISPR screen of \citet{Norman2019} (Section~\ref{sec:results-transfer}); no
dataset-specific parameters are tuned.

\subsection{Reachability as membership in a convex cone}
\label{sec:methods-cone}
Two properties of loss-of-function perturbations determine the geometry. First, a CRISPRi
knockdown removes rather than adds function, so each atom $\Emat_{:,j}$ enters a combination
with a fixed orientation. Second, applying a set of knockdowns together produces, to first
order, the sum of their individual effects; there is no natural notion of a
``negative amount'' of a knockdown. The set of cell-state shifts reachable by a non-negative
combination of knockdowns is therefore the \emph{non-negative conic hull} of the measured
effect vectors,
\begin{equation}
\cone \;=\; \Bigl\{\, \Emat\wvec \;:\; \wvec \in \Rnn^{P} \,\Bigr\}
\;\subseteq\; \mathbb{R}^{G},
\label{eq:cone}
\end{equation}
a closed convex cone. Deciding whether a target $\dvec$ is reachable is deciding whether
$\dvec\in\cone$ (up to a non-negative scale, since only direction is claimed). We quantify
proximity by projecting $\dvec$ onto $\cone$: the reachable component is the solution of the
non-negative least-squares (NNLS) program
\begin{equation}
\wvec^{\star} \;=\; \arg\min_{\wvec \,\geq\, 0}\; \bigl\lVert \Emat\wvec - \dvec \bigr\rVert_2^{2},
\qquad
\hat{\dvec} \;=\; \Emat\wvec^{\star},
\label{eq:nnls}
\end{equation}
solved with the Lawson--Hanson active-set algorithm \citep{LawsonHanson1974}. The
\emph{reachable cosine} $\cos(\hat{\dvec},\dvec)$ measures how much of the target's direction
the cone can express, and the residual $\dvec-\hat{\dvec}$ is the part it cannot. A target is
declared reachable when its reachable cosine is both high in absolute terms and far above the
null of Section~\ref{sec:methods-null}; the same solve returns the support of $\wvec^{\star}$,
whose largest entries constitute a \emph{minimal ranked knockdown recipe}. Ordering genes by
a greedy forward selection that adds at each step the knockdown giving the largest marginal
gain in reachable cosine yields the reachability spectrum reported throughout.

\subsection{Infeasibility certificates via convex duality}
\label{sec:methods-cert}
The value of a feasibility oracle is that ``no'' is provable, not merely observed. When
$\dvec\notin\cone$, the separating-hyperplane theorem for convex cones (a form of Farkas'
lemma \citep{Farkas1902}) guarantees a vector $\bm{y}$ with $\Emat^{\top}\bm{y}\le 0$ and
$\dvec^{\top}\bm{y}>0$: a hyperplane with the whole cone on one side and the target strictly
on the other. This $\bm{y}$ is exactly the negative gradient of the NNLS objective at the
optimum, so it is produced by the same solve rather than a separate computation. The
Karush--Kuhn--Tucker (KKT) stationarity conditions of \eqref{eq:nnls} provide the numerical
proof: at the optimum, active (positive-weight) atoms have zero reduced gradient and inactive
atoms have non-positive reduced gradient. We report the maximum stationarity violation as the
certificate residual; across all reported solves it is at machine precision
($\approx\!1.1\times10^{-11}$), confirming true optimality rather than early stopping. The
biological reading of the certificate is constructive: the positive coordinates of the
residual direction---the part of the target orthogonal to everything the cone can
supply---name genes the target requires to be \emph{raised} that no knockdown combination can
deliver. These become falsifiable CRISPRa (activation) hypotheses, turning an infeasibility
result into a redirected experiment rather than a null one.

\subsection{Signed loss-/gain-of-function decomposition}
\label{sec:methods-decomp}
To separate ``cannot be reached because it needs activation'' from ``cannot be reached for
other reasons,'' we decompose the target into three orthogonal budgets as fractions of its
squared norm: the component captured by the non-negative knockdown cone
(loss-of-function, LOF); the component lying along directions that require increased
expression (gain-of-function, GOF), obtained by projecting the residual onto the
positive-expression orthant; and an irreducible remainder aligned with neither
(``neither''). We report this in two conventions. The \emph{signed} decomposition partitions
the target by the sign of each coordinate's required change---an orthogonal split along
mutually polar cones in the sense of \citet{Moreau1962}; the \emph{one-sided}
decomposition instead attributes the residual after the non-negative projection. For
T\textsubscript{H}2$\rightarrow$T\textsubscript{H}1 these give LOF/GOF/neither of
$39\%/25\%/35\%$ (signed) and $39\%/31\%/30\%$ (one-sided). We report both because they
answer subtly different questions and neither is privileged; the qualitative conclusion---%
that loss-of-function alone accounts for well under half of the shift---is common to both.

\subsection{Held-out validation and null models}
\label{sec:methods-null}
Two null models guard against two distinct failure modes, each addressing a different question.

\paragraph{Held-out-gene generalization.} An in-sample reachable cosine can be inflated by a
solver freely reweighting thousands of atoms to fit a fixed target; it measures expressivity
of the dictionary, not predictive validity. We therefore report a \emph{held-out-gene}
cosine: we fit the NNLS weights using only a random subset of measured genes (rows of
$\Emat$ and $\dvec$) and evaluate the reachable cosine on the held-out genes, so the recipe
must generalize to expression coordinates it was not fit on. For
T\textsubscript{H}2$\rightarrow$T\textsubscript{H}1 the held-out cosine is $0.448$, against an
in-sample value of $0.627$; the gap is the cost of generalization, and we quote the
held-out number as the headline throughout.

\paragraph{Shuffled-target null.} To calibrate what reachable cosine arises by chance from a
high-dimensional cone, we recompute the held-out reachability against $60$ random targets
generated by shuffling the gene labels of $\dvec$, preserving its marginal magnitude
distribution while destroying its biological structure. For
T\textsubscript{H}2$\rightarrow$T\textsubscript{H}1 the null has mean $-0.005$, standard
deviation $0.019$, and maximum $0.029$; the observed held-out cosine sits at $z\approx24$
above it. This null answers ``is the target special,'' distinct from held-out-gene
generalization, which answers ``does the recipe transfer to unseen coordinates.''

\paragraph{Analytic anisotropy null.} Because a shuffled-target null over an atlas of targets
is expensive, we derive a closed-form approximation to the expected null cosine as a function
of two geometric quantities---the alignment of the target with the dominant direction of the
cone and the effective anisotropy of the atom cloud (Supplementary
Fig.~\ref{fig:supp-matrix} legend). Validated against the shuffle null on both synthetic and
real CD4\textsuperscript{+} data, it reproduces the empirical null to Pearson $r=0.995$
(in-sample) and $0.998$ (held-out), with maximum absolute error $0.047$ and $0.07$
respectively, replacing on the order of a thousand target shuffles with ten to twenty NNLS
solves. This makes the atlas-wide reliability analysis of Section~\ref{sec:results-atlas}
tractable while keeping the empirical shuffle null as the reference for the headline result.

\paragraph{Positive controls.} We verify that the geometry recovers known biology using
canonical T-helper regulators as controls, \emph{not} as discoveries. On the
toward-T\textsubscript{H}1 track, GATA3---the master T\textsubscript{H}2 transcription factor,
which must be removed---is reached at rank $155/6871$ ($97.7$th percentile, cosine $+0.101$)
on the loss-of-function side, while TBX21---the master T\textsubscript{H}1 factor, which must
be raised---is correctly anti-aligned under knockdown at rank $6775/6871$ ($1.4$th
percentile). Because these regulators are already reported by the source screen
\citep{Zhu2025}, their recovery validates the operator without being claimed as a novel
finding.

\subsection{Robustness and reinforcement analyses}
\label{sec:methods-robustness}
Five analyses probe whether the headline verdict is an artifact of a specific modelling or
evaluation choice. All reuse \texttt{reachability.py} without modification.

\paragraph{Non-negativity ablation.} To test whether the biological non-negativity constraint
earns its place, we compare the held-out cosine of the constrained NNLS solve against an
unconstrained least-squares solve (which permits negative weights, i.e. treating a knockdown
as if it could be run in reverse) and against the single best-aligned generator, across all
$12$ atlas cells. We report the cosine cost of the constraint, the number of biologically
unrealizable negative weights the unconstrained solution requires, and the fraction of the
unconstrained cosine the cone recovers.

\paragraph{Reachable-cosine ceiling.} Because a knockdown-only cone cannot express
gain-of-function directions, the achievable held-out cosine has a modality-imposed ceiling. We
compute this ceiling as $\sqrt{\smash[b]{f_{\mathrm{LOF}}}}$, where $f_{\mathrm{LOF}}$ is the
loss-of-function fraction of the signed decomposition, and express each headline cosine as a
fraction of its own ceiling so that a ``modest'' absolute number is reported against what is
attainable in principle rather than against $1$.

\paragraph{DEG-magnitude weighting.} Uniform-gene metrics can dilute a real signal by
averaging over thousands of unchanged genes, a critique raised for perturbation-response
evaluation by \citet{Mejia2025}. We recompute every reachable cosine and held-out $z$ under a
weighting $w_g=\lvert d_g\rvert$ that concentrates the metric on the genes the target actually
moves, and we recalibrate the dynamic-range fraction against a shuffled-gene floor and an
interpolated-duplicate ceiling under the same weighting. The unweighted metric remains the
default and reproduces the published numbers exactly; the weighting is a stress test, not a
re-definition.

\paragraph{Cross-cell-type transfer.} To ask whether the operator generalizes beyond a single
cell type, we apply it to the genome-scale essential-gene Perturb-seq screens of
\citet{Replogle2022} in K562 and RPE1 cells ($843$ shared single-gene perturbations over
$2{,}832$ shared genes). We evaluate transfer at three levels: whether a single
perturbation's effect \emph{direction} agrees across cell types (matched cosine against a
shuffled-gene null); whether the reachability \emph{verdict} agrees (graded by the residual
norm within- versus cross-basis, since a binary verdict is uninformative when an $843$-atom
cone is highly expressive); and whether the minimal \emph{recipe} agrees (Jaccard overlap of
selected atoms against a random-selection null).

\paragraph{Certificate stability and external cross-reference.} We test the activation
certificate two ways. For reproducibility, we refit the certificate on a random half of the
gene axis and measure the rank agreement of the top genes against the full-data ranking and
across independent halves (Spearman $\rho$ over $12$ random splits, against a label-shuffle
null). For biological interpretation, we retrospectively cross-reference the top certificate
genes against the literature direction-of-effect and against Open Targets genetic-association
scores for the relevant T\textsubscript{H}1-driven autoimmune diseases. This cross-reference is
an exploratory annotation, not a pre-registered test, and we report it as such.

\subsection{Software and reproducibility}
\label{sec:methods-software}
The oracle is implemented in a single module (\texttt{reachability.py}) built on standard
NNLS routines; all reported quantities derive from the published effect matrices of
\citet{Zhu2025} and \citet{Norman2019} with no additional model fitting. Every numeric value
in this manuscript is drawn from a locked facts sheet tied to the analysis outputs, and each
figure is generated directly from those outputs; figures were checked for exact
correspondence to the facts sheet before inclusion.

